\subsection*{Tiling further: two cooperative groups sharing operands}

Section~\ref{sec:3-4} introduced \texttt{parallel p1 by 2 :\allowbreak{} group-\allowbreak{}4} --- with MMA, both groups can share the same B tile while loading different rows of A.

\begin{lstlisting}[language=croqtile,caption={Tiling further: two cooperative groups sharing operands}]
#define TILE_M 128
#define TILE_N 128
#define TILE_K 128
#define WARP_M 64
#define WARP_N 64
#define WARP_K 16
__co__ void matmul_wgmma_2group(global f16 [M, K] lhs, global f16 [N, K] rhs, global f16 [M, N] output) {
  parallel {block_m, block_n} by [cdiv(M, TILE_M), cdiv(N, WARP_N)] : block {
    shared f16 [TILE_M, TILE_K] lhs_load_s;
    shared f16 [WARP_N, TILE_K] rhs_load_s;
    shared f16 [TILE_M, WARP_N] output_s;
    mc = mma.fill.f32 0.0f;
    foreach {iv_k} in [cdiv(K, TILE_K)] {
      dma.copy lhs.subspan(TILE_M, TILE_K).at(block_m, iv_k) => lhs_load_s;
      dma.copy rhs.chunkat(block_n, iv_k) => rhs_load_s;
      parallel p1 by 2 : group-4 {
        foreach {iv_warp} in [cdiv(TILE_K, WARP_K)] {
          ma = mma.load lhs_load_s.subspan(WARP_M, TILE_K).at(p1, 0).chunkat(_, iv_warp);
          mb = mma.load rhs_load_s.chunkat(_, iv_warp);
          mma.row.row mc, ma, mb;
        }
      }
    }
    parallel p1 by 2 : group-4
      mma.store mc, output_s.subspan(WARP_M, WARP_N).at(p1, 0);
    dma.copy output_s => output.subspan(TILE_M, WARP_N).at(block_m, block_n);
  }
}
\end{lstlisting}

With \co@TILE_M = 128@ and a 64-row group tile, \co@p1@ selects either the upper or lower half of the block's output rows. The groups load different rows of A and share \co@rhs_load_s@, since both need the same B columns and K slice.

The four-step syntax stays invariant; only the parallel decomposition of the tile changes.

\begin{table}[htbp]
\centering\small
\begin{tabularx}{\linewidth}{@{}lXXX@{}}
\toprule
\textbf{Aspect} & One group & One group-4 & Two group-4 groups \\
\midrule
Operand sharing & N/A & N/A & B tile shared, A rows split by \texttt{p1} \\
Threads & 32 & 128 & 256 \\
Accumulator & \texttt{mma.fill 0.0} & \texttt{mma.\allowbreak{}fill.\allowbreak{}f16 0.\allowbreak{}0f} & \texttt{mma.\allowbreak{}fill.\allowbreak{}f32 0.\allowbreak{}0f} \\
\bottomrule
\end{tabularx}
\end{table}

The two groups must store disjoint output row ranges. They may read the same right-hand tile because those reads do not modify it, but each group needs its own accumulator for its assigned rows. This separates operand reuse from output ownership: sharing an input does not imply sharing a partial sum.
